\section{Doxygen-dokumentasjon}
Her presenteres en detaljert dokumentasjon av heissystemets kodestruktur i Doxygen-format.

\subsection{Datastrukturer}
Heissystemet benytter følgende hovedstrukturer:

\begin{lstlisting}[style=cpp, caption=Heissystemets tilstander]
/**
 * @brief Mulige tilstander for heisen
 */
typedef enum { 
    DOOR_OPEN = 0,       /**< Døren er åpen */
    STANDING_STILL = 1,  /**< Heisen står stille med lukket dør */
    DRIVING_UP = 2,      /**< Heisen beveger seg oppover */
    DRIVING_DOWN = 3     /**< Heisen beveger seg nedover */
} State;

/**
 * @brief Heisens retning
 */
typedef enum { 
    CURRENT_DIR_UP = 0,    /**< Siste retning var oppover */
    CURRENT_DIR_DOWN = 1,  /**< Siste retning var nedover */
} Direction;

/**
 * @brief Datastruktur som inneholder heissystemets tilstand
 */
typedef struct state_data {
    State state;          /**< Nåværende tilstand */
    int btnStates[4][3];  /**< Knappetrykk for alle etasjer og knappetyper */
    Direction dir;        /**< Nåværende/siste retning */
} state_data;
\end{lstlisting}

\subsection{Tilstandsmodul (states.h)}
Modulen inneholder funksjoner for å håndtere de ulike tilstandene i heisen.

\begin{lstlisting}[style=cpp, caption=Hovedfunksjoner i tilstandsmodulen]
/**
 * @brief Håndterer tilstanden hvor heisen står stille
 * 
 * Når heisen er i denne tilstanden, sjekker den etter
 * nye bestillinger og bestemmer neste tilstand.
 * 
 * @param data Peker til heisens tilstandsdata
 * @param floor Nåværende etasje (-1 hvis mellom etasjer)
 */
void standing_still(state_data *data, int floor);

/**
 * @brief Åpner heisdøren
 * 
 * Aktiverer døråpningslampen og venter i 3 sekunder
 * før den slår av lampen.
 * 
 * @param data Peker til heisens tilstandsdata
 */
void open_door(state_data *data);

/**
 * @brief Håndterer tilstanden med åpen dør
 * 
 * Styrer tidtakeren for åpen dør og går tilbake til standing_still
 * når tidtakeren utløper. Sjekker også for hindring og stoppknapp.
 * 
 * @param data Peker til heisens tilstandsdata
 * @param floor Nåværende etasje (-1 hvis mellom etasjer)
 */
void door_open_state(state_data *data, int floor);

/**
 * @brief Håndterer kjøretilstanden (opp eller ned)
 * 
 * Kontrollerer heisbevegelsen og sjekker om den skal stoppe
 * ved etasjer. Håndterer stoppknappen og overganger til andre tilstander.
 * 
 * @param data Peker til heisens tilstandsdata
 * @param floor Nåværende etasje (-1 hvis mellom etasjer)
 * @param dir Nåværende bevegelsesretning
 */
void driving(state_data *data, int floor, MotorDirection dir);

/**
 * @brief Nullstiller alle bestillinger og slår av alle knappelys
 * 
 * Brukes når stoppknappen trykkes eller i andre situasjoner
 * hvor alle bestillinger skal kanselleres.
 * 
 * @param data Peker til heisens tilstandsdata
 */
void clear_all_orders(state_data *data);

/**
 * @brief Sjekker om det er en hindring i døren
 * 
 * Håndterer hindring ved å stoppe heisen, åpne døren,
 * og nullstille alle bestillinger.
 * 
 * @param data Peker til heisens tilstandsdata
 * @param floor Nåværende etasje (-1 hvis mellom etasjer)
 * @return 1 hvis hindring er aktiv, 0 ellers
 */
int check_obstruction(state_data *data, int floor);
\end{lstlisting}

\subsection{Inputmodul (input.h)}
Modulen håndterer knappetrykk og lyssignaler.

\begin{lstlisting}[style=cpp, caption=Hovedfunksjoner i inputmodulen]
/**
 * @brief Sjekker for knappetrykk og oppdaterer tilstandsdata
 * 
 * Skanner alle heisknapper (hall opp, hall ned, cab) på alle etasjer
 * og oppdaterer tilstandsdata og knappelys.
 * Ignorerer knappetrykk hvis stoppknappen er trykket eller hindring er aktiv.
 * 
 * @param data Peker til heisens tilstandsdata
 */
void checkForInput(state_data *data);

/**
 * @brief Slår av alle knappelys for en spesifikk etasje
 * 
 * Brukes når en etasje betjenes for å nullstille bestillinger og slå av lys.
 * 
 * @param data Peker til heisens tilstandsdata
 * @param floor Etasjen hvor knappelys skal slås av
 */
void turnOffButtonLamps(state_data *data, int floor);
\end{lstlisting}

\subsection{Oppstartsmodul (startup.h)}
Modulen håndterer initialisering av heisen.

\begin{lstlisting}[style=cpp, caption=Hovedfunksjoner i oppstartsmodulen]
/**
 * @brief Initialiserer heisen ved oppstart ved å kjøre til en definert posisjon
 * 
 * Ignorerer alle bestillinger og stoppknappen under initialisering.
 * Kjører heisen til nærmeste etasje hvis den er mellom etasjer.
 * 
 * @return 1 når initialiseringen er fullført
 */
int elevator_startup();

/**
 * @brief Slår av alle knappelys
 * 
 * Nullstiller alle knappelys ved systemoppstart.
 */
void clear_buttons();
\end{lstlisting}

\subsection{Grensesnitt mot heismaskinvare (elevio.h)}
Modulen inneholder funksjoner for kommunikasjon med heismaskinvaren.

\begin{lstlisting}[style=cpp, caption=Hovedfunksjoner i maskinvaregrensesnittet]
/**
 * @brief Initialiserer heismaskinvaren
 * 
 * Må kalles en gang før andre funksjoner benyttes.
 */
void elevio_init(void);

/**
 * @brief Setter motorens retning
 * 
 * @param dirn Ønsket retning (DIRN_UP, DIRN_DOWN, DIRN_STOP)
 */
void elevio_motorDirection(MotorDirection dirn);

/**
 * @brief Setter lampen for en spesifikk knapp
 * 
 * @param floor Etasjen (0-indeksert)
 * @param button Knappetypen
 * @param value 1 for å slå på, 0 for å slå av
 */
void elevio_buttonLamp(int floor, ButtonType button, int value);

/**
 * @brief Setter etasjeindikatoren
 * 
 * @param floor Etasjen som skal vises (0-indeksert)
 */
void elevio_floorIndicator(int floor);

/**
 * @brief Kontrollerer døråpningslampen
 * 
 * @param value 1 for å slå på, 0 for å slå av
 */
void elevio_doorOpenLamp(int value);

/**
 * @brief Kontrollerer stoppknapplampen
 * 
 * @param value 1 for å slå på, 0 for å slå av
 */
void elevio_stopLamp(int value);

/**
 * @brief Leser statusen til en ringeknapp
 * 
 * @param floor Etasjen (0-indeksert)
 * @param button Knappetypen
 * @return 1 hvis knappen er trykket, 0 ellers
 */
int elevio_callButton(int floor, ButtonType button);

/**
 * @brief Leser etasjesensoren
 * 
 * @return Nåværende etasje (0-indeksert), eller -1 hvis mellom etasjer
 */
int elevio_floorSensor(void);

/**
 * @brief Leser status på stoppknappen
 * 
 * @return 1 hvis stoppknappen er trykket, 0 ellers
 */
int elevio_stopButton(void);

/**
 * @brief Leser status på hindringsbryteren
 * 
 * @return 1 hvis hindring er aktiv, 0 ellers
 */
int elevio_obstruction(void);
\end{lstlisting}

\subsection{Hovedprogram (main.c)}
Hovedprogrammet for heiskontrollsystemet.

\begin{lstlisting}[style=cpp, caption=Hovedfunksjonen]
/**
 * @file main.c
 * @brief Hovedprogram for heiskontrollsystemet
 *
 * Denne filen inneholder hovedløkken for heiskontrollsystemet.
 * Den implementerer en tilstandsmaskin som håndterer de ulike
 * heistilstandene og overganger mellom dem basert på inndata og hendelser.
 *
 * @author Erik Bohne & Lars Lien
 */

/**
 * @brief Hovedfunksjon for heiskontrollsystemet
 *
 * Initialiserer heismaskinvaren, utfører oppstartskalibrering,
 * og kjører hovedkontrollløkken som implementerer heistilstandsmaskinen.
 *
 * @return 0 ved normal avslutning, -1 ved feil
 */
int main(){
    // Initialiser maskinvare
    elevio_init();
    
    // Initialiser tilstandsdatastruktur
    state_data sData = {STANDING_STILL, {{0,0,0},{0,0,0},{0,0,0},{0,0,0}}, CURRENT_DIR_DOWN};
    
    // Utfør heisoppstartskalibrering
    elevator_startup();

    // Hovedkontrollløkke
    while(1){
        // Oppdater etasjesensoravlesning
        floor = elevio_floorSensor();
        if(floor >= 0){
            elevio_floorIndicator(floor);
        }
        
        // Sjekk først for hindringer - dette har høyest prioritet
        if (check_obstruction(&sData, floor)) {
            continue;
        }
        
        // Sjekk for knappeinndata
        checkForInput(&sData);

        // Tilstandsmaskinimplementasjon
        switch (sData.state) {
        case STANDING_STILL:
            standing_still(&sData, floor);
            break;
        case DRIVING_UP:
            driving(&sData, floor, DIRN_UP);
            break;
        case DRIVING_DOWN:
            driving(&sData, floor, DIRN_DOWN);
            break;
        case DOOR_OPEN:
            door_open_state(&sData, floor);
            break;
        default:
            return -1;
            break;
        }
    }
    
    return 0;
}
\end{lstlisting}